% HMT provenance: output/EXTREMA_MEDIA_RAZON_NARRACION_Y_REVISION_20260916/
% ARTICULO/ES/source/libro/01_acoplamiento.tex, lib01:cadena-registro and lib01:K.
% Direction: ARTICULO/ES/source/sections/k_direccion_dimensional.tex.
% Status of the networks, charts, and flows in this appendix: FORMALIZACION_NUEVA.
% The extraction of K and the projector P_3 are RESULTADO_RECUPERADO.
\section{Planar boundary measurement, Plücker coordinates, and positive mutations}
\label{app:network-cluster}

The dodecaphasic publication of the enriched state admits a positive
realization whose network, coordinates, and chart changes can be written
out completely. The starting composition is
\[
 \mathrm{APP}\longrightarrow\mathrm{TRIT}\longrightarrow\mathrm{TPK}
 \longrightarrow x_{\rm term}\longrightarrow\mathcal R_{12}
 \longrightarrow(b^{90},b^{120},Q_{\rm TPK})
 \longrightarrow U_{\rm sgn}\longrightarrow K.
\]
This reading acts on the enriched state preserving the two sheets,
orientation, carry, memory, and boundary. The already generated register is
\[
 K=(234,543,140,729,659,824,621,58,914,794,146,601),
 \qquad Q=\sum_{r=1}^{12}K_r=6263.
\]
Its extraction is recovered exactly from the two channels and total charge:
\[
 b^{90}=(-495,-116,-684,108,601,-90,-173,-88,313,560,-397,461),
\]
\[
 b^{120}=(-425,-281,-481,671,-255,30,475,-543,680,251,6,-128).
\]
On the position orbits \((r,r+3,r+6,r+9)\), for \(r=1,2,3\),
let \(B_r=\sum_{j=0}^3b^{120}_{r+3j}\) and
\(d_{r,j}=b^{90}_{r+3j}\). We obtain
\[
 s_1=(Q+2B_1+B_2)/3,\quad s_2=s_1-B_1,\quad s_3=s_2-B_2,
\]
\[
 U^{(r)}=(s_r,d_{r,1}-d_{r,3},d_{r,0}+d_{r,2},d_{r,0}-d_{r,2}),
 \qquad K|_r=\tfrac14H_4U^{(r)},
\]
\[
 H_4=\begin{pmatrix}1&1&1&1\\1&1&-1&-1\\1&-1&1&-1\\1&-1&-1&1\end{pmatrix},
 \qquad H_4^2=4I_4.
\]
The network constructed below realizes the positive values produced by
this extraction. Its combinatorial form is a subsequent formalization
of the register, with stated origin and orientation.

\subsection{Directed acyclic network and determinants of nonintersecting paths}

For a positive register of length \(L\), set
\[
 d_r=K_r/Q>0,\qquad t_0=0,\qquad
 t_j=\sum_{r=1}^jd_r\quad(1\le j\le L).
\]
Thus \(t_L=1\). The \(L\) separations produce \(L+1\) marked
points; in the dodecaphasic case there are thirteen points.
The matrix realizing them is
\begin{equation}
 C(d)=\begin{pmatrix}1&1&\cdots&1\\t_0&t_1&\cdots&t_L\end{pmatrix}.
 \label{eq:nc-path-matrix}
\end{equation}
The Grassmannian \(\operatorname{Gr}(2,L+1)\) identifies rank-two
matrices under invertible changes of the two rows. Its positive part
consists of classes admitting all ordered minors of order two strictly
positive.

Construct a directed network with vertices \(a_j,b_j\), for
\(0\le j\le L\), and sinks \(q_j\). There are horizontal edges
\(a_{j-1}\to a_j\) and \(b_{j-1}\to b_j\), of weight one;
vertical edges \(a_j\to b_j\), of weight \(d_j\), for \(j\ge1\);
and edges \(b_j\to q_j\), of weight one. The ordered sources are
\(b_0,a_0\). The two rows are drawn horizontally, the vertical edges
descend, and the sinks exit below the lower row. This drawing is planar,
and the orientation is acyclic. The boundary measurement \(M_{ij}\)
is the sum of the products of weights of all paths from source \(i\)
to sink \(q_j\).
The path network has two sources and \(L+1\) sinks, hence \(L+3\)
external connections. The matrix it realizes has \(L+1\) columns.
The plabic seed introduced below has precisely \(L+1\) boundary
vertices and constitutes another presentation of its coordinates;
these two boundary counts are kept explicit.

\begin{theorem}[Boundary measurement and register minors]
\label{thm:nc-network}
The boundary measurement of this network is \(M=C(d)\).
For \(0\le i<j\le L\),
\begin{equation}
 \Delta_{ij}(C(d))=t_j-t_i=\sum_{r=i+1}^{j}d_r>0.
 \label{eq:nc-minor-gap}
\end{equation}
The same number is the sum of weights of the vertex-disjoint path
pairs \(b_0\to q_i\), \(a_0\to q_j\).
\end{theorem}
\begin{proof}
The lower source has a unique path to each sink, of weight one.
A path from \(a_0\) to \(q_j\) descends exactly once, along an
edge of index \(1\le r\le j\); its weight is \(d_r\). This
proves the matrix formula. The lower path to \(q_i\) occupies all
vertices \(b_0,\ldots,b_i\). A path from the other source to
\(q_j\) is therefore disjoint from it precisely when it descends
at \(i<r\le j\). Its sum of weights is the right-hand side of
\eqref{eq:nc-minor-gap}. The pair with exchanged destinations always
shares a vertex. Expanding the determinant cancels all intersecting
pairs and retains exactly the disjoint ones. Here the determinantal
mechanism of nonintersecting paths has been verified directly for the
specific network.
\end{proof}

Identification of planar-network measurements with cells of the
nonnegative Grassmannian provides their subsequent geometric recognition;
see A. Postnikov,
\href{https://arxiv.org/abs/math/0609764}{\emph{Total positivity,
Grassmannians, and networks}}. The preceding proof also gives a finite
formula for every minor without using a metrological value.

The family \(d_r>0\), \(\sum_rd_r=1\), has dimension \(L-1\).
Its map to the Grassmannian is injective: if the rows of \(C(d)\)
and \(C(d')\) generate the same plane, the second rows differ by an
affine transformation; the conditions \(t_0=t'_0=0\) and
\(t_L=t'_L=1\) fix that transformation. In particular, for \(L=12\)
one obtains an eleven-dimensional family inside the twenty-two-dimensional
positive cell. The point lies in the top cell; the family retains its
own dimension. The thirteen points and twelve separations thus have
different mathematical roles.

Allowing \(d_r\ge0\) preserves nonnegativity. When \(d_r=0\),
consecutive columns \(r-1\) and \(r\) coincide and form parallel
elements of the matroid. The rank remains two because \(t_0=0\)
and \(t_L=1\). The bases are exactly the pairs \(i<j\) crossing
some positive separation. This degeneration produces strata of
consecutive incidence. A zero column, which would give a loop,
corresponds to a different operation from the coincidence of two
nonzero columns.

\subsection{Triangulation seed and exchange relation}

Let \(N=L+1\), and consider the vertices \(0,\ldots,L\) of a
convex polygon in cyclic order. To each side or diagonal \((i,j)\),
assign the coordinate \(A_{ij}=\Delta_{ij}(C(d))\), with ordered
indices. The fan triangulation from vertex zero has diagonals
\((0,j)\), \(2\le j\le L-1\). Its sides are frozen variables,
and its diagonals are mutable variables. The quiver is obtained by
cyclically orienting the three sides of each triangle and canceling
opposite arrow pairs. This is a seed of type \(A_{N-3}\); for the
dodecaphasic register, its type is \(A_{10}\).

The seed has an explicit plabic realization. Place a black vertex inside
each triangle; at the polygon vertices, place a white vertex if it is
incident to a diagonal, and a black one otherwise. Join the centers to
the three corners of their triangles, remove the original sides, and
contract edges between same-colored vertices. Add a boundary leaf at
each marked vertex. This construction follows the triangulation
algorithm of Kodama and Williams,
\href{https://people.math.harvard.edu/~williams/papers/KW-ArxivVersion.pdf}
{\emph{KP solitons and total positivity for the Grassmannian}, Algorithm 12.1}.
The coordinate labels of this seed are evaluated at the minors already
produced by the two-row network. For a mutable face, the
\(\mathcal X\)-type coordinate is the ratio of products of incident
\(\mathcal A\) coordinates prescribed by its quiver column; in a
quadrilateral it will be written explicitly as a cross-ratio.
Changing the graph changes the chart of these minors while preserving
their provenance.

\begin{theorem}[Positive exchange in the register realization]
\label{thm:nc-ptolemy}
For four vertices \(a<b<c<d\),
\begin{equation}
 A_{ac}A_{bd}=A_{ab}A_{cd}+A_{ad}A_{bc}.
 \label{eq:nc-ptolemy}
\end{equation}
Consequently, the diagonal change \((a,c)\leftrightarrow(b,d)\)
is evaluated through
\[
 A_{bd}=\frac{A_{ab}A_{cd}+A_{ad}A_{bc}}{A_{ac}}>0.
\]
All triangulations provide compatible positive charts on the same
realization \(C(d)\).
\end{theorem}
\begin{proof}
Substituting \(A_{ij}=t_j-t_i\) reduces the equality to
\[
 (t_c-t_a)(t_d-t_b)
 =(t_b-t_a)(t_d-t_c)+(t_d-t_a)(t_c-t_b).
\]
Expanding both sides proves the identity. The differences are positive;
therefore each flip preserves a positive evaluation, and the quotient
is defined. Connectivity of triangulations by flips allows passage
between them through these exchanges. The cluster interpretation of
the Plücker relation is the subsequent recognition developed by
J. Scott in
\href{https://arxiv.org/abs/math/0311148}{\emph{Grassmannians and Cluster
Algebras}}. The algebraic proof of the specific evaluation is already
contained in the displayed identity.
\end{proof}

\subsection{Refinement of separations and naturality of exchanges}

Let \(B\ge2\). At depth \(n\), each separation \(d_r\) is
divided into \(B^n\) equal consecutive separations:
\[
 d^{(n)}_{(r-1)B^n+c}=d_r/B^n,
 \qquad 1\le c\le B^n.
\]
This defines \(LB^n+1\) points and a matrix \(C_n\). If
\(\iota_n(j)=Bj\) identifies a point with its same-position
descendant, then
\begin{equation}
 C_{n+1}|_{\iota_n(\{0,\ldots,LB^n\})}=C_n,
 \qquad
 \Delta_{\iota_n(i),\iota_n(j)}(C_{n+1})=\Delta_{ij}(C_n).
 \label{eq:nc-refinement}
\end{equation}
The refined network replaces each vertical edge of weight \(d\)
by a sequence of \(B\) descent positions, with weights \(d/B\).
The path sum between inherited endpoints remains \(d\).
Composition of refinements preserves this equality because finite
sums are grouped associatively.

\begin{proposition}[Naturality on inherited endpoints]
\label{prop:nc-inherited}
Every exchange \eqref{eq:nc-ptolemy} among four inherited endpoints
commutes exactly with refinement. Every old triangulation extends to
one of the refined polygon while preserving its interior triangles:
the polygons added between consecutive endpoints are triangulated
outside those triangles. Flips between old diagonals then preserve
their exchange relations.
\end{proposition}
\begin{proof}
The first conclusion follows by substituting the equality in
\eqref{eq:nc-refinement} for each minor. For the second, preserve
the chords joining consecutive old endpoints and triangulate each
outer pocket formed by the new marks. These regions have disjoint
interiors. The two triangles incident to an old diagonal remain
unchanged, and so do their quadrilateral and Ptolemy identity.
\end{proof}

This naturality specifies the endpoints on which it acts. An old
side becomes a diagonal of the larger polygon when new vertices
are added along its boundary arc. A formerly frozen variable can
therefore be mutable in the enlarged problem. The preceding
compatibility is exact compatibility of the evaluation system and
inherited flips. An inclusion of seed patterns with the same
coefficient regime is obtained by keeping these chords and the new
auxiliary diagonals frozen. This states the operation that preserves
the seed; increasing depth and performing a mutation remain different
operations.

\subsection{Projective diagonal action and invariance of cross-ratios}

If \(\lambda_j>0\), the column action
\(C\mapsto C\operatorname{diag}(\lambda_0,\ldots,\lambda_L)\) gives
\[
 A_{ij}\longmapsto\lambda_i\lambda_j A_{ij}.
\]
It preserves positivity and all strata defined by vanishing minors.
The cross-ratio associated with a quadrilateral is
\begin{equation}
 X_{abcd}=\frac{A_{ab}A_{cd}}{A_{ad}A_{bc}}>0.
 \label{eq:nc-crossratio}
\end{equation}
The four factors \(\lambda_a\lambda_b\lambda_c\lambda_d\) cancel
exactly. Consequently, every positive diagonal column action fixes
\(X_{abcd}\), even if it depends on a parameter and has determinant
one. Changing column weights preserves the projective point of each
column; their cross-ratios express this preservation. For the register
to change the projective shape, it must act on the ordered separations.

\subsection{Register-directed positive deformation of separations}

The corpus fixes the representation of the twelve positions and its
orthogonal projector \(P_3\) onto the three-dimensional sector.
With \(P_{11}=I_{12}-J_{12}/12\), the recovered direction is
\[
 u=P_3P_{11}K=P_3K\in\mathbf1^\perp.
\]
In the stated chart, setting \(r=\sqrt5\), its exact evaluation is
\begin{align*}
u=(&-495/4-233r/10,\ 403/4+169r/10,\ -403/4-169r/10,\\
   &495/4+233r/10,\ 601/4+1031r/10,\ 203/4+211r/5,\\
   &-203/4-211r/5,\ -601/4-1031r/10,\ 313/4+569r/10,\\
   &162+219r/20,\ -162-219r/20,\ -313/4-569r/10).
\end{align*}
We retain \(\sum_ru_r=0\) and
\(\|u\|^2=(6638585+2275584\sqrt5)/20>0\).

For \(s\in\mathbb R\), define the deformation of the weights by
\begin{equation}
 d_r(s)=\frac{K_re^{su_r}}{\sum_{\ell=1}^{12}K_\ell e^{su_\ell}},
 \qquad t_j(s)=\sum_{r=1}^jd_r(s),\qquad C(s)=C(d(s)).
 \label{eq:nc-gap-flow}
\end{equation}
The formula uses the already generated register and its direction.
It is a new analytic realization of the reader; the parameter \(s\)
measures its deformation. When the internal output
\(\rho_A=\pi_{\rm HMT}/\varphi_{\rm HMT}^2\) is used, writing
\(s=\tau\log\rho_A\) expresses exactly the same family.
This reparametrization acts after the generation of the constants.
Identifying \(s\) or \(\tau\) with a physical time requires the
temporal law of the corresponding reader.

\begin{theorem}[Shape change, positivity, and compatibility with depth]
\label{thm:nc-shape-flow}
The family \eqref{eq:nc-gap-flow} belongs to
\(\operatorname{Gr}_{>0}(2,13)\) for every real \(s\) and preserves
\(t_0=0\), \(t_{12}=1\). For \(i<j\), let
\[
 \mu_{ij}(s)=
 \frac{\sum_{r=i+1}^jK_ru_re^{su_r}}{\sum_{r=i+1}^jK_re^{su_r}},
 \qquad
 \mu(s)=\frac{\sum_rK_ru_re^{su_r}}{\sum_rK_re^{su_r}}.
\]
Then
\begin{equation}
 \frac{d}{ds}\log A_{ij}(s)=\mu_{ij}(s)-\mu(s),
 \quad
 \frac{d}{ds}\log X_{abcd}(s)
 =\mu_{ab}(s)+\mu_{cd}(s)-\mu_{ad}(s)-\mu_{bc}(s).
 \label{eq:nc-crossratio-derivative}
\end{equation}
The deformation commutes with every refinement assigning the children
of separation \(r\) the weight \(d_r/B\) and the same velocity \(u_r\).
\end{theorem}
\begin{proof}
All exponential factors and all \(K_r\) are positive. Normalization
makes the separations sum to one; the minors are sums of positive
separations by \eqref{eq:nc-minor-gap}. Differentiating the logarithm of
\[
 A_{ij}(s)=\frac{\sum_{r=i+1}^jK_re^{su_r}}
                      {\sum_rK_re^{su_r}}
\]
gives the first identity. The combination of four logarithms in
\eqref{eq:nc-crossratio} cancels the four global means with their
signs and gives the second. Finally, replacing \(K_r\) by \(B\)
copies of \(K_r/B\) with identical \(u_r\) preserves numerator
and denominator when the children are grouped. The equality holds
for all parameters, at any depth, and for every composition of refinements.
\end{proof}

There is an exact witness that this family changes shape. For the
first four endpoints,
\begin{equation}
 X_{0123}(0)=\frac{1560}{23711},\qquad
 \left.\frac{d}{ds}\log X_{0123}(s)\right|_{s=0}
 =-\frac{309899}{917}-\frac{268596}{4585}\sqrt5<0.
 \label{eq:nc-nontrivial-witness}
\end{equation}
The calculation uses only the first three separations
\((234,543,140)\) and their already recovered velocities. The strict
inequality proves an effective projective motion. The Ptolemy identities
remain exact during this motion, since they apply to the columns
\((1,t_j(s))\) at every parameter value. The register thus controls
a positive shape deformation and, simultaneously, a system of charts
that remains compatible with it.

The proved scope brings together a rank-two planar network, its explicit
measurement, a triangulation plabic seed, all its mutations by flips,
naturality of inherited endpoints, and a deformation of cross-ratios
directed by \(K\). Each operation has a verifiable domain and result.
Comparison of this realization with enriched temporal dynamics and
the material branch uses their respective reading maps; no coordinate
identity replaces those maps.
